% GENERATED FILE. Edit canonical lessons, then run npm run generate:books.
% canonical-source-hash: fnv1a64:afd8480e4c507754
% canonical-lessons: ES-C578-desacuerdo, ES-C578-discrepar, ES-C578-oponerse, ES-C578-partidario, ES-C578-contrario

\chapter{No estoy de acuerdo}
\label{ch:es-no-estoy-de-acuerdo}
\hlchaptermodality{578}

\begin{chapteropening}
\textbf{By the end of this chapter:} \emph{I can say politely that I disagree, that I am against something or in favour of it.}

Complete the last lesson of chapter 578: I can say politely that I disagree, that I am against something or in favour of it.
\end{chapteropening}

\section[el desacuerdo]{el desacuerdo — the disagreement}

\label{lesson:ES-C578-desacuerdo}



\begin{quote}
Before the new one: say the Spanish for reasonable, then the Spanish for logical.
\end{quote}

\subsection*{You'll want to know: el desacuerdo}
\textbf{el desacuerdo} — \textquotedblleft{}the disagreement\textquotedblright{}.

\begin{quote}
Estoy en desacuerdo con esa idea.
\end{quote}

\begin{quote}
\emph{I disagree with that idea.}
\end{quote}

\emph{Estar en desacuerdo con} is the opposite of \emph{estar de acuerdo con}.

\begin{cousinweb}
\emph{desacuerdo} is built on \emph{acuerdo} with the prefix \emph{des-}; \emph{acuerdo} comes from \emph{acordar}, inherited from spoken Latin \emph{accordare}, built on \emph{cor} 'heart' (EtymDB).
\end{cousinweb}

\begin{grammarlens}[title={What you've built}]
One more word for telling, explaining and describing at B1.
\end{grammarlens}

\subsection*{Your turn}
\begin{itemize}
\raggedright
  \item \emph{Say it:} \emph{el desacuerdo}
  \item \emph{Say it:} \emph{el desacuerdo}, once more
  \item \emph{Say it:} a sentence of your own with \emph{el desacuerdo}
  \item \emph{Recall:} say the Spanish for reasonable, then the Spanish for logical, then say \emph{el desacuerdo} again
\end{itemize}

\begin{tcolorbox}[breakable,colback=teal!4,colframe=teal!35!black,title={Before you move on}]
What does el desacuerdo mean? (\textquotedblleft{}The disagreement\textquotedblright{}.) Say it once more.
\end{tcolorbox}
\section[discrepar]{discrepar — to disagree, to differ}

\label{lesson:ES-C578-discrepar}



\begin{quote}
Before the new one: say the Spanish for logical, then the Spanish for the disagreement.
\end{quote}

\subsection*{You'll want to know: discrepar}
\textbf{discrepar} — \textquotedblleft{}to disagree, to differ\textquotedblright{}.

\begin{quote}
Discrepo un poco de tu opinión.
\end{quote}

\begin{quote}
\emph{I disagree a little with your opinion.}
\end{quote}

A polite, formal verb: \emph{discrepar de} someone or something.

\begin{cousinweb}
\emph{discrepar} is a learned borrowing from Latin \emph{discrepare} 'to be out of tune, to differ', built from \emph{dis-} 'apart' and \emph{crepare} 'to rattle, to creak' (EtymDB).
\end{cousinweb}

\begin{grammarlens}[title={What you've built}]
One more word for telling, explaining and describing at B1.
\end{grammarlens}

\subsection*{Your turn}
\begin{itemize}
\raggedright
  \item \emph{Say it:} \emph{discrepar}
  \item \emph{Say it:} \emph{discrepar}, once more
  \item \emph{Say it:} a sentence of your own with \emph{discrepar}
  \item \emph{Recall:} say the Spanish for logical, then the Spanish for the disagreement, then say \emph{discrepar} again
\end{itemize}

\begin{tcolorbox}[breakable,colback=teal!4,colframe=teal!35!black,title={Before you move on}]
What does discrepar mean? (\textquotedblleft{}To disagree, to differ\textquotedblright{}.) Say it once more.
\end{tcolorbox}
\section[oponerse]{oponerse — to oppose, to be against}

\label{lesson:ES-C578-oponerse}



\begin{quote}
Before the new one: say the Spanish for the disagreement, then the Spanish for to disagree.
\end{quote}

\subsection*{You'll want to know: oponerse}
\textbf{oponerse} — \textquotedblleft{}to oppose, to be against\textquotedblright{}.

\begin{quote}
Me opongo a ese cambio.
\end{quote}

\begin{quote}
\emph{I am against that change.}
\end{quote}

Like \emph{poner}: \emph{me opongo}. Always with \emph{a}: \emph{oponerse a algo}.

\begin{cousinweb}
\emph{oponerse} is a learned borrowing from Latin \emph{opponere} 'to set against', \emph{ob} + \emph{ponere} 'to put', rebuilt on \emph{poner}; EtymDB does not list it, but the Latin gloss is in Lewis \& Short.
\end{cousinweb}

\begin{grammarlens}[title={What you've built}]
One more word for telling, explaining and describing at B1.
\end{grammarlens}

\subsection*{Your turn}
\begin{itemize}
\raggedright
  \item \emph{Say it:} \emph{oponerse}
  \item \emph{Say it:} \emph{oponerse}, once more
  \item \emph{Say it:} a sentence of your own with \emph{oponerse}
  \item \emph{Recall:} say the Spanish for the disagreement, then the Spanish for to disagree, then say \emph{oponerse} again
\end{itemize}

\begin{tcolorbox}[breakable,colback=teal!4,colframe=teal!35!black,title={Before you move on}]
What does oponerse mean? (\textquotedblleft{}To oppose, to be against\textquotedblright{}.) Say it once more.
\end{tcolorbox}
\section[partidario]{partidario — in favour; supporter}

\label{lesson:ES-C578-partidario}



\begin{quote}
Before the new one: say the Spanish for to disagree, then the Spanish for to oppose.
\end{quote}

\subsection*{You'll want to know: partidario}
\textbf{partidario} — \textquotedblleft{}in favour; supporter\textquotedblright{}.

\begin{quote}
Soy partidario de trabajar en casa.
\end{quote}

\begin{quote}
\emph{I am in favour of working from home.}
\end{quote}

\emph{Ser partidario de} + noun or infinitive. As a noun: \emph{los partidarios del plan}.

\begin{cousinweb}
\emph{partidario} is built on \emph{partido} 'a side, a party', from \emph{partir} 'to divide', which is inherited from Latin \emph{partire}, built on \emph{pars} 'part' (Lewis \& Short).
\end{cousinweb}

\begin{grammarlens}[title={What you've built}]
One more word for telling, explaining and describing at B1.
\end{grammarlens}

\subsection*{Your turn}
\begin{itemize}
\raggedright
  \item \emph{Say it:} \emph{partidario}
  \item \emph{Say it:} \emph{partidario}, once more
  \item \emph{Say it:} a sentence of your own with \emph{partidario}
  \item \emph{Recall:} say the Spanish for to disagree, then the Spanish for to oppose, then say \emph{partidario} again
\end{itemize}

\begin{tcolorbox}[breakable,colback=teal!4,colframe=teal!35!black,title={Before you move on}]
What does partidario mean? (\textquotedblleft{}In favour; supporter\textquotedblright{}.) Say it once more.
\end{tcolorbox}
\section[contrario]{contrario — opposed, against; opposite}

\label{lesson:ES-C578-contrario}



\begin{quote}
Before the new one: say the Spanish for to oppose, then the Spanish for in favour.
\end{quote}

\subsection*{You'll want to know: contrario}
\textbf{contrario} — \textquotedblleft{}opposed, against; opposite\textquotedblright{}.

\begin{quote}
Mi hermano es contrario a la idea.
\end{quote}

\begin{quote}
\emph{My brother is against the idea.}
\end{quote}

\emph{Al contrario} — on the contrary. \emph{Lo contrario} — the opposite.

\begin{cousinweb}
\emph{contrario} is a learned borrowing from Latin \emph{contrarius} 'opposite', built on \emph{contra} 'against' (Lewis \& Short).
\end{cousinweb}

\begin{grammarlens}[title={What you've built}]
One more word for telling, explaining and describing at B1.
\end{grammarlens}

\subsection*{Your turn}
\begin{itemize}
\raggedright
  \item \emph{Say it:} \emph{contrario}
  \item \emph{Say it:} \emph{contrario}, once more
  \item \emph{Say it:} a sentence of your own with \emph{contrario}
  \item \emph{Recall:} say the Spanish for to oppose, then the Spanish for in favour, then say \emph{contrario} again
\end{itemize}

\begin{tcolorbox}[breakable,colback=teal!4,colframe=teal!35!black,title={Before you move on}]
What does contrario mean? (\textquotedblleft{}Opposed, against; opposite\textquotedblright{}.) Say it once more.
\end{tcolorbox}
